\scanpage{116}
% line head | text | 65
\mspage{65}
% line 1–2 | text | 1.3. Flèches dans Cat
\textbf{1.3.}\quad \deleted{\unsure{Foncteurs}} Flèches dans $\mathrm{Cat}$, de source ou but un objet
recollé $(X,U,T)$.

% line 3–4 | text | Quand (X,U,T) est donnée
Quand $(X,U,T)$ est donnée par \unsure{leur|leurs} donnée de
recollement
% line 5 | display | U, T, H
\[
U,T,\qquad H:U^{\mathrm o}\times T\longrightarrow\mathrm{Ens},
\]
% line 6 | text | les flèches
les \deleted{foncteurs} flèches
% line 7 | display | f, X prime, X
\[
f:X'\longrightarrow X
\]
% line 8–9 | text | dans Cat se décrivent
dans $\mathrm{Cat}$ se décrivent immédiatement \unsure{ainsi} : \deleted{\unsure{pour}}
la donnée de $f$ définit d'abord
% line 10 | display | U prime and T prime
\[
U'=\unsure{j|f}^{-1}(U),\qquad T'=f^{-1}(T)
\]
% line 11–13 | text | couple d'un ouvert
couple d'un ouvert de $X'$ et du fermé complémentaire,
\deleted{\unsure{et}} de sorte que cette \deleted{\unsure{d.}} donnée équivaut
à celle de

% line 14 | text | 1°) un couple
$\underline{1}^{\circ})$ un couple $(U',T')$ dans $X'$

% line 15–17 | text | 2°) un morphisme
$2^{\circ})$ un morphisme (de $\unsure{\mathscr T|\mathcal T}$) de $(X',U',T')$
\unsure{dans} $(X,U,T)$, ce qui équivaut par 1.1.
à un \unsure{triple}
% line 18–20 | display | f on U prime, f on T prime, u
\[
\left\{\begin{aligned}
&f_{U'}:U'\to U\\
&f_{T'}:T'\unsure{-|\to}T\\
&\unsure{u|v}:H_{U',T'}\longrightarrow H\circ(f_{U'}^{\mathrm o}\times f_{T'})
\end{aligned}\right.
\]
% line 21 | text | où on a posé
où on a posé
% line 22–23 | display | H with subscripts, its value at u prime and t prime
\[
\left\{\begin{aligned}
&H_{U',T'}:{U'}^{\mathrm o}\times T'\longrightarrow\mathrm{Ens}\\
&H_{U',T'}(u',t')=\operatorname{Hom}_{X'}(u',t'),
\end{aligned}\right.
\]
% line 24–25 | text | i.e. H is the gluing datum
i.e. $H_{U',T'}$ est la donnée de recollement
qui correspond au triple $(X',U^{\unsure{\prime\prime|\prime}},T')$ de $\mathscr T$.
% archive sheet number: 116
